\documentclass[12pt]{article}

\usepackage[a4paper, margin=1in]{geometry}
\usepackage{amsmath,amssymb,amsthm}
\usepackage{mathtools}
\usepackage{mathrsfs}
\usepackage{enumitem}
\usepackage{xcolor}
\usepackage[hidelinks]{hyperref}
\usepackage{fontspec}

\defaultfontfeatures{Ligatures=TeX}
\IfFontExistsTF{Times New Roman}
  {\setmainfont{Times New Roman}}
  {\setmainfont{TeX Gyre Termes}}
\IfFontExistsTF{Menlo}
  {\setmonofont{Menlo}}
  {\setmonofont{Latin Modern Mono}}

\newtheorem{theorem}{Theorem}[section]
\newtheorem{proposition}[theorem]{Proposition}
\newtheorem{lemma}[theorem]{Lemma}
\newtheorem{corollary}[theorem]{Corollary}
\theoremstyle{definition}
\newtheorem{definition}[theorem]{Definition}
\theoremstyle{remark}
\newtheorem{remark}[theorem]{Remark}

\newcommand{\tightlist}{%
  \setlength{\itemsep}{0pt}\setlength{\parskip}{0pt}}

\title{%
  \textbf{Foundation II: Shared Manifestation}\\
  \textbf{A Conditional Two-Source Program for Shared Accessible Sectors}%
}

\author{
  Sidong Liu, PhD \\
  \small iBioStratix Ltd \\
  \small \texttt{sidongliu@hotmail.com}
}

\date{April 2026}

\begin{document}
\emergencystretch=2em
\raggedbottom
\hbadness=5000
\vbadness=5000

\maketitle

\begin{abstract}
Foundation~I of the Consciousness Series~\cite{FoundationI} constructs a
conditional chain from a single source datum
\(S_{\mathrm{source}}=(\mathcal H_U,\mathcal A_U,\Omega,\{\mathcal E_t\})\)
to an admissible accessible pair \((\mathcal A_c,\omega)\), which then
feeds into the inseparability theorem of Paper~I~\cite{PaperI}. That
chain is single-source: it produces one observer's manifest sector. The
present paper, Foundation~II, asks the next structural question: when two
source chains share the same ambient substrate, under what conditions
does a non-trivial \emph{shared} manifest sector emerge?

We define a multi-source datum as a family of source data sharing a
common ambient triple \((\mathcal H_U,\mathcal A_U,\Omega)\), each
equipped with its own decoherence semigroup. No global tensor-product
decomposition is assumed; the multi-source structure is encoded entirely
through subalgebra inclusions, following the non-canonicality of
tensor-product factorizations in the Type~III\(_1\) setting established
in~\cite{HAFFA}. We then introduce joint-stability conditions under
which a candidate shared sector \(\mathcal A_{\mathrm{sh}}\) is
simultaneously stable for both sources.

The paper's main result is a conditional two-source foothold, not a
general multi-source closure theorem. Under named joint-stability and
compatible-projection conditions, we construct a
shared stable hull
\(\Phi_{\mathrm{sh}}^{\mathrm{st}}(\mathfrak F_{\mathrm{sh}})\) from
the overlap of the two foreground frames, prove that it is a von Neumann
subalgebra carrying a faithful restricted state, and show that, under an
additional admissibility input, the resulting pair
\((\Phi_{\mathrm{sh}}^{\mathrm{st}}(\mathfrak F_{\mathrm{sh}}),
\omega_{\mathrm{sh}}^{\mathrm{st}})\) lies in the theorem domain of
Paper~I. The
shared sector therefore inherits the
geometry--modular-flow inseparability of the single-source case.

After the main theorem, we analyze a regime-dependent common/private
structure: under a shared/private regime assumption, the shared sector
is the locus on the overlap algebra where the two manifest states agree,
while observer-dependent discrepancy can only appear outside that shared
sector. Finally, we state a conditional \(N\)-source extension under
pairwise-threshold and exchangeability assumptions, and record the open
problems that remain, including the quantitative relationship between
consensus strength and geometric stability.

This paper treats only the mathematical bridge from multi-source data to
shared accessible sectors. It does not claim a general \(N\)-source
theorem, does not provide a quantitative consensus--stability bound, and
does not derive the candidate shared sector from bare multi-source
dynamics alone. It also does not attempt to settle questions of
intersubjective ontology.
\end{abstract}

\clearpage
\tableofcontents
\clearpage

\section{Introduction}\label{sec:intro}

Foundation~I of the Consciousness Series~\cite{FoundationI} addressed
the single-source problem. Starting from a source datum
\[
S_{\mathrm{source}}
=
(\mathcal H_U,\mathcal A_U,\Omega,\{\mathcal E_t\}_{t\ge 0}),
\]
it isolated a conditional chain
\[
S_{\mathrm{source}}
\to
\mathcal O_A
\to
\mathcal O_B
\to
(\mathcal A_c,\omega),
\]
where \(\mathcal O_A\) is a persistent projected record process,
\(\mathcal O_B\) is a budget-forced foreground frame, and
\((\mathcal A_c,\omega)\) is the manifest accessible pair that then
feeds into the geometry--modular-flow inseparability theorem of
Paper~I~\cite{PaperI}. That was the correct scope for the first
foundation paper: it explained how one observer-facing manifest sector
could arise from one source chain.

The present paper asks the next structural question. If two source
chains share the same ambient substrate, under what conditions does a
non-trivial \emph{shared} manifest sector arise? This is the minimum
multi-source problem behind any later discussion of a shared world. The
paper does not attempt to settle the full intersubjective problem in one
step. It does not claim a general \(N\)-source theorem, does not derive
a quantitative consensus-strength/stability law, and does not aim to
settle questions of ontology. Its narrower aim is to identify a
theorem-level foothold for \emph{shared manifestation existence}.

That calibration matters. In the current programme there are two
separate bridges beyond the already-written Consciousness Trilogy
results. The first concerns whether entangled or correlated source
chains can support a common manifest sector at all. The second concerns
how the strength of that commonality controls the quantitative stability
of the resulting geometry. Foundation~II addresses only the first. The
second remains a separate Bridge~3 problem and should not be smuggled
into the present paper under vague language about ``consensus'' or
``robustness''.

The paper therefore adopts the smallest honest target. Its main theorem
is proved in the two-source case. The theorem does not directly state
that a full shared geometry has been derived. Its first conclusion is
weaker and mathematically safer: under named compatibility and
joint-stability inputs, one can construct a
\emph{shared admissible manifest sector}
\[
(\mathcal A_{\mathrm{sh}},\omega_{\mathrm{sh}})
\subseteq
(\mathcal A_c^{(1)},\omega_1)\cap(\mathcal A_c^{(2)},\omega_2),
\]
which is stable under the relevant source-side dynamics and can then be
fed into the already-existing Paper~I machinery. Geometry enters only as
a corollary once that admissible shared sector is in hand.
The main result should therefore be read as a conditional inclusion
theorem and theorem-level foothold, not as a full emergence theorem for
shared worlds.

At the heuristic level, this shared-sector question is related to the
quantum Darwinism and spectrum-broadcast-structure literature, where the
objective part of a quantum description is identified with redundantly
accessible pointer information carried across multiple environment
fragments~\cite{Zurek2003,Zurek2009,BrandaoPianiHorodecki2015,KorbiczHorodeckiHorodecki2014}.
The present paper does not translate those frameworks into a theorem
about Type~III\(_1\) ambient algebras. It only records the
operator-algebraic analogue needed here: a shared foreground overlap
together with a shared stable hull.

Three design choices govern the whole paper.
\begin{enumerate}
\item \textbf{Two sources first.} The paper starts at \(N=2\), because
  the two-source case is the minimum non-trivial setting in which one
  can ask whether a shared manifest sector exists at all. General
  \(N\)-source statements are postponed to a conditional extension
  section.
\item \textbf{No tensor-product starting point.} The multi-source datum
  is formulated inside a common ambient algebra rather than from a
  canonical tensor decomposition. This follows the HAFF viewpoint that
  physically relevant subsystem structure is selected by stability and
  accessibility rather than fixed in advance~\cite{HAFFA,HAFFD}.
\item \textbf{Common/private decomposition is secondary.} After the main
  theorem, we analyze how each observer's full manifest algebra may be
  written, under a regime assumption, as a shared sector plus a private
  complement. That decomposition is treated as a structural consequence,
  not as a competing main theorem.
\end{enumerate}

The paper's mathematical contributions are correspondingly limited and
explicit.
\begin{enumerate}
\item We define a \emph{multi-source datum} as a common ambient
  Type~III\(_1\) standard-form algebra together with a family of
  source-specific decoherence semigroups.
\item We separate the new inputs into three named
  joint-stability conditions: shared-sector invariance,
  shared-state existence with faithfulness, and projection/kernel
  compatibility.
\item In the two-source case, we define the shared foreground overlap
  and the associated shared stable hull, and prove the abstract
  minimality properties of that hull.
\item Under an explicit candidate shared-sector input, we prove an
  inclusion theorem showing that the shared stable hull lies inside the
  common manifest carrier, and therefore yields a shared admissible
  manifest sector.
\item Under a shared/private regime assumption, we analyze the resulting
  common/private structure by identifying the shared sector as the locus
  of agreement on the overlap algebra and by localizing overlap-level
  discrepancy outside that sector.
\end{enumerate}

The role of conditional language is exactly the same as in
Foundation~I. Wherever a piece can be internalized cleanly, we
internalize it. Wherever it cannot, we isolate it and keep it explicit.
The present theorem does not derive the existence of a candidate shared
sector from bare multi-source dynamics or entanglement alone. It starts
from such a sector as named input and analyzes the stable closure that
follows from it.
The main open points here are not hidden: model-independent joint
stability, compatible Nakajima--Zwanzig projections, and the
well-definedness of common/private decomposition in the Type~III\(_1\)
setting.

The paper is organized as follows. Section~\ref{sec:datum} defines the
multi-source datum and the joint-stability package. Section~\ref{sec:two}
states and proves the two-source main theorem. Section~\ref{sec:cp}
analyzes the common/private structure under a regime assumption.
Section~\ref{sec:nsource} gives a conditional \(N\)-source extension.
Section~\ref{sec:open} records the open problems. Section~\ref{sec:ledger}
collects the named conditional inputs. Appendix~\ref{app:premises}
summarizes the imported Genesis and Consciousness results used by the
paper.

\section{Multi-Source Datum and Joint-Stability Inputs}\label{sec:datum}

The first task is to define the source-side object language without
overloading it with later conditional interface data.

\begin{definition}[Multi-source datum]
A \emph{multi-source datum} is a tuple
\[
S_{\mathrm{multi}}
:=
(\mathcal H_U,\mathcal A_U,\Omega,\{\mathcal E_t^{(i)}\}_{i=1}^N),
\]
where:
\begin{enumerate}
\item \((\mathcal H_U,\mathcal A_U,\Omega)\) is a Type~III\(_1\) von
  Neumann algebra in standard form, with \(\Omega\) cyclic and
  separating for \(\mathcal A_U\);
\item for each \(i\in\{1,\dots,N\}\),
  \(\{\mathcal E_t^{(i)}\}_{t\ge 0}\) is a pointwise
  \(\sigma\)-weakly continuous semigroup of normal completely positive
  unital maps on \(\mathcal A_U\).
\end{enumerate}
\end{definition}

\begin{remark}[Heisenberg-picture convention]
The semigroups \(\{\mathcal E_t^{(i)}\}\) act on observables, so the
paper uses the Heisenberg-picture convention throughout. The continuity
assumption in Definition~2.1 is therefore pointwise \(\sigma\)-weak
continuity on \(\mathcal A_U\), not norm continuity.
\end{remark}

\begin{remark}[A deliberately thin datum]
The multi-source datum contains only the common ambient standard-form
triple and the family of source-side dynamical semigroups. It does
\emph{not} include Nakajima--Zwanzig projections, foreground frames,
accessible algebras, or any private/shared decomposition. Those are not
primitive data; they are interface or output objects that may or may not
be obtainable from the datum under additional assumptions.
\end{remark}

\begin{remark}[Why no tensor-product factorization]\label{rem:notensor}
The present paper does not assume a global tensor-product decomposition
\[
\mathcal H_U \cong \bigotimes_{i=1}^N \mathcal H_i
\qquad \text{or} \qquad
\mathcal A_U \cong \bar\otimes_{i=1}^N \mathcal A_i.
\]
This is deliberate. In the HAFF programme, the physically relevant
subalgebras are selected by accessibility and stability rather than by a
canonical prior factorization~\cite{HAFFA,HAFFD}. In the Type~III
setting, treating the subsystem split as primitive would import exactly
the kind of structure the Genesis programme is designed to derive rather
than presuppose.
\end{remark}

The single-source outputs from Foundation~I are now imported only as
downstream conditional objects.

\begin{definition}[Individual manifest outputs]
Let \(S_{\mathrm{multi}}\) be a multi-source datum. For each source
index \(i\), an \emph{individual manifest output} means a pair
\[
(\mathcal A_c^{(i)},\omega_i)
\]
obtained from the Foundation~I chain for the \(i\)-th source, where
\(\mathcal A_c^{(i)}\subseteq \mathcal A_U\) is a von Neumann
subalgebra and \(\omega_i\) is a faithful normal state on
\(\mathcal A_c^{(i)}\).
\end{definition}

\begin{remark}
Nothing in Definition~2.3 claims that every source in a multi-source
datum automatically produces an individual manifest output. That is
already the conditional content of Foundation~I~\cite{FoundationI}.
Foundation~II takes those single-source outputs as antecedent material
when the two-source main theorem is stated.
\end{remark}

\begin{definition}[Candidate shared sector]
Given individual manifest outputs
\((\mathcal A_c^{(i)},\omega_i)_{i=1}^N\), a
\emph{candidate shared sector} is a von Neumann subalgebra
\[
\mathcal A_{\mathrm{sh}}^{\mathrm{cand}}
\subseteq
\bigcap_{i=1}^N \mathcal A_c^{(i)}.
\]
\end{definition}

The new source-side content begins with the joint-stability package. We
state it in three layers so that later failures can be localized
cleanly.

\begin{definition}[JS-A: shared-sector invariance and semigroup compatibility]
\label{def:JSA}
Let \(S_{\mathrm{multi}}\) be a multi-source datum and let
\(\mathcal A_{\mathrm{sh}}^{\mathrm{cand}}\subseteq\mathcal A_U\) be a
candidate shared sector. We say that
\((\mathcal A_{\mathrm{sh}}^{\mathrm{cand}},\{\mathcal E_t^{(i)}\})\)
satisfies \emph{JS-A} if:
\begin{enumerate}
\item \(\mathcal E_t^{(i)}(A)\in \mathcal A_{\mathrm{sh}}^{\mathrm{cand}}\)
  for every \(A\in\mathcal A_{\mathrm{sh}}^{\mathrm{cand}}\), every
  \(i\), and every \(t\ge 0\);
\item the ambient modular flow
  \(\sigma_s^{\omega,\mathrm{tot}}\) of
  \(\omega(\cdot)=\langle\Omega,\cdot\,\Omega\rangle\) satisfies
  \[
  \sigma_s^{\omega,\mathrm{tot}}
  (\mathcal A_{\mathrm{sh}}^{\mathrm{cand}})
  =
  \mathcal A_{\mathrm{sh}}^{\mathrm{cand}}
  \qquad
  \forall s\in\mathbb R;
  \]
\item the restricted semigroups commute on the candidate shared sector:
  \[
  \mathcal E_t^{(i)}\mathcal E_u^{(j)}(A)
  =
  \mathcal E_u^{(j)}\mathcal E_t^{(i)}(A)
  \qquad
  \forall A\in \mathcal A_{\mathrm{sh}}^{\mathrm{cand}},
  \ \forall t,u\ge 0.
  \]
\end{enumerate}
\end{definition}

\begin{remark}[Status of the JS-A commutativity clause]
The exact commutativity requirement in Definition~\ref{def:JSA} is a
strong regime assumption. The present paper does not claim that it
follows model-independently from the bare multi-source datum. A weaker
variant would replace strict commutativity by asymptotic commutativity
in the long-time limit, or by commutativity only on the shared
fixed-point algebra. The stronger form is used here because it makes the
shared stable-carrier construction transparent; the weaker variants are
left open.
\end{remark}

\begin{definition}[JS-B: shared-state existence and faithfulness]
\label{def:JSB}
Let \(\mathcal A_{\mathrm{sh}}^{\mathrm{cand}}\subseteq \mathcal A_U\)
be a candidate shared sector. We say that \(\mathcal A_{\mathrm{sh}}^{\mathrm{cand}}\)
satisfies \emph{JS-B} if there exists a faithful normal state
\[
\omega_{\mathrm{sh}}
\in
(\mathcal A_{\mathrm{sh}}^{\mathrm{cand}})_*
\]
such that:
\begin{enumerate}
\item \(\omega_{\mathrm{sh}}\circ\mathcal E_t^{(i)}=\omega_{\mathrm{sh}}\)
  on \(\mathcal A_{\mathrm{sh}}^{\mathrm{cand}}\) for every
  \(i\) and every \(t\ge 0\);
\item \(\omega_{\mathrm{sh}}\) is compatible with the individual
  manifest states in the sense that
  \[
  \omega_i|_{\mathcal A_{\mathrm{sh}}^{\mathrm{cand}}}
  =
  \omega_{\mathrm{sh}}
  \qquad
  \forall i\in\{1,\dots,N\}.
  \]
\end{enumerate}
\end{definition}

\begin{lemma}[Vector-state simplification of JS-B]\label{lem:vectorJSB}
Assume the common ambient vector-state regime:
\[
\omega(\cdot)=\langle\Omega,\cdot\,\Omega\rangle
\qquad
\text{on } \mathcal A_U.
\]
If \(\mathcal M\subseteq \mathcal A_U\) is any von Neumann subalgebra,
then the restricted state
\[
\omega|_{\mathcal M}
\]
is faithful on \(\mathcal M\).
\end{lemma}

\begin{proof}
Because \(\mathcal M\subseteq \mathcal A_U\subseteq B(\mathcal H_U)\),
both algebras act on the same Hilbert space \(\mathcal H_U\). Since
\(\Omega\) is separating for \(\mathcal A_U\), the set
\(\mathcal A_U'\Omega\) is dense in \(\mathcal H_U\). The inclusion
\(\mathcal M\subseteq \mathcal A_U\) implies \(\mathcal A_U'\subseteq
\mathcal M'\), hence \(\mathcal M'\Omega\supseteq \mathcal A_U'\Omega\)
is dense. Therefore \(\Omega\) is separating for \(\mathcal M\).
Hence if \(A\in\mathcal M_+\) and
\(\omega(A)=\langle\Omega,A\Omega\rangle=0\), then
\(A^{1/2}\Omega=0\), so \(A^{1/2}=0\) and therefore \(A=0\).
\end{proof}

\begin{remark}[Why JS-B is still kept explicit]
Lemma~\ref{lem:vectorJSB} shows that in the strict common-vector-state
regime, faithfulness on a shared subalgebra is automatic. We keep JS-B
explicit nevertheless, for two reasons. First, later variants may allow
source-dependent faithful normal states not all induced by the same
ambient vector. Second, even in the common-vector regime, one still
wants a named place where state compatibility across sources is recorded
explicitly rather than silently assumed.
\end{remark}

\begin{remark}[Strength of the JS-B stationarity clause]
The stationarity requirement in Definition~\ref{def:JSB} is strong. It
amounts to asking that the shared state lie in the common fixed-state
sector of all source semigroups when restricted to
\(\mathcal A_{\mathrm{sh}}^{\mathrm{cand}}\). The present paper does not
claim that such a common stationary state is generic. A weaker variant
would replace exact stationarity by long-time convergence or approximate
stationarity on the shared sector; those relaxations are left open.
\end{remark}

\begin{definition}[JS-C: compatible persistent projections]\label{def:JSC}
Let \(S_{\mathrm{multi}}\) be a multi-source datum, and fix a candidate
shared sector \(\mathcal A_{\mathrm{sh}}^{\mathrm{cand}}\). A pair of
Nakajima--Zwanzig projections
\[
\mathcal P^{(1)},\mathcal P^{(2)}:\mathcal A_U\to\mathcal A_U
\]
is called \emph{\(\mathcal A_{\mathrm{sh}}^{\mathrm{cand}}\)-compatible}
if:
\begin{enumerate}
\item each \(\mathcal P^{(i)}\) is a normal idempotent projection used
  in the Stage~A output of the \(i\)-th source chain;
\item they commute,
  \[
  \mathcal P^{(1)}\mathcal P^{(2)}
  =
  \mathcal P^{(2)}\mathcal P^{(1)}
  =:\mathcal P_{\mathrm{sh}};
  \]
\item \(\mathcal P_{\mathrm{sh}}(\mathcal A_U)\) is a von Neumann
  subalgebra containing
  \(\mathcal A_{\mathrm{sh}}^{\mathrm{cand}}\), and
  \(\mathcal P_{\mathrm{sh}}|_{\mathcal A_{\mathrm{sh}}^{\mathrm{cand}}}
  = \mathrm{id}\);
\item the projected predual dynamics induced by the pair admit a common
  dual memory kernel
  \[
  \mathcal K_{\mathrm{sh},*}(t,s)
  \]
  on the shared projected sector, meaning that there exists a single
  kernel governing the shared projected dynamics on that sector; the
  definition does not require literal equality of the two individual
  source kernels before projection to the shared sector.
\end{enumerate}
\end{definition}

\begin{remark}[Status of JS-C]
Definition~\ref{def:JSC} is not imported from a standard open-systems
textbook. It is a new named compatibility condition introduced for the
present paper. The point is not to pretend that a model-independent
notion of compatible Nakajima--Zwanzig projections is already standard.
The point is to isolate exactly what later proofs need from such a
notion. The background projection formalism is still the standard
Nakajima--Zwanzig one~\cite{Nakajima1958,Zwanzig1960}; what is new here
is only the multi-source compatibility layer. A weaker variant would
require only that the long-time asymptotic behaviour of the two
projected kernels agree on the shared sector, rather than demanding a
common kernel at all times. The present paper uses the stronger form
because it simplifies the Stage~A construction; the weaker variant is
left as an open relaxation.
\end{remark}

\section{The Two-Source Shared Manifest Sector}\label{sec:two}

We now specialize to \(N=2\). The paper's main theorem is entirely
two-source.

\subsection{Joint source interface}

Let
\[
S_{\mathrm{multi}}^{(2)}
=
(\mathcal H_U,\mathcal A_U,\Omega,\{\mathcal E_t^{(1)}\},\{\mathcal E_t^{(2)}\})
\]
be a two-source datum. Assume that each source separately satisfies the
named Foundation~I hypotheses required to produce a Stage~A output,
Stage~B output, and manifest pair
\[
(\mathcal A_c^{(i)},\omega_i),
\qquad
i=1,2.
\]

\begin{lemma}[Joint Stage~A output from compatible persistent projections]
\label{lem:jointA}
Let \(S_{\mathrm{multi}}^{(2)}\) be a two-source datum. Assume there
exists a candidate shared sector
\(\mathcal A_{\mathrm{sh}}^{\mathrm{cand}}\subseteq
\mathcal A_c^{(1)}\cap\mathcal A_c^{(2)}\) satisfying JS-C. Then the
commuting product
\[
\mathcal P_{\mathrm{sh}}
:=
\mathcal P^{(1)}\mathcal P^{(2)}
\]
defines a shared projected record object
\[
\mathcal O_A^{\mathrm{sh}}
:=
(\mathcal P_{\mathrm{sh}},\omega_{\mathrm{sh},t},\mathcal K_{\mathrm{sh},*}(t,s)),
\]
where \(\omega_{\mathrm{sh},t}\) is the projected normal state on the
shared sector and \(\mathcal K_{\mathrm{sh},*}(t,s)\) is the common dual
memory kernel from Definition~\ref{def:JSC}.
\end{lemma}

\begin{proof}
Commutativity of \(\mathcal P^{(1)}\) and \(\mathcal P^{(2)}\) implies
that \(\mathcal P_{\mathrm{sh}}\) is again an idempotent projection.
Because both \(\mathcal P^{(1)}\) and \(\mathcal P^{(2)}\) are normal
and they commute, their composition is again normal. By JS-C,
\(\mathcal P_{\mathrm{sh}}(\mathcal A_U)\) is a von Neumann subalgebra
containing \(\mathcal A_{\mathrm{sh}}^{\mathrm{cand}}\), and the common
dual kernel \(\mathcal K_{\mathrm{sh},*}(t,s)\) is part of the named
input. This is exactly the shared Stage~A object required for later
stages.
\end{proof}

\subsection{Joint foreground overlap}

For each source \(i=1,2\), let
\(\mathfrak F^{(i)}\) denote the foreground frame produced by the
Foundation~I budget-forced compression step, with foreground realization
\[
\iota_i:V_{\mathrm{fg}}(\mathfrak F^{(i)})\hookrightarrow
\mathcal A_U^{\mathrm{sa}}.
\]

\begin{definition}[Shared foreground overlap]
\label{def:sharedfg}
The \emph{shared foreground overlap} is the operator-side intersection
\[
V_{\mathrm{fg}}(\mathfrak F_{\mathrm{sh}})
:=
\iota_1\bigl(V_{\mathrm{fg}}(\mathfrak F^{(1)})\bigr)
\cap
\iota_2\bigl(V_{\mathrm{fg}}(\mathfrak F^{(2)})\bigr)
\subseteq
\mathcal A_U^{\mathrm{sa}}.
\]
Whenever this intersection is non-zero and finite-dimensional, we choose
a basis
\[
O_1^{\mathrm{sh}},\dots,O_{k_{\mathrm{sh}}}^{\mathrm{sh}}
\in
V_{\mathrm{fg}}(\mathfrak F_{\mathrm{sh}})
\]
and write
\[
\mathcal A_{\mathfrak F_{\mathrm{sh}}}
:=
W^*(1,O_1^{\mathrm{sh}},\dots,O_{k_{\mathrm{sh}}}^{\mathrm{sh}}).
\]
\end{definition}

\begin{proposition}[Dimension bound for the shared foreground]
\label{prop:sharedfgdim}
Assume \(\mathfrak F^{(1)}\) and \(\mathfrak F^{(2)}\) have effective
dimensions \(k_1^*\) and \(k_2^*\), respectively, and that the shared
foreground overlap of Definition~\ref{def:sharedfg} is finite
dimensional. Then
\[
k_{\mathrm{sh}}
\le
\min\{k_1^*,k_2^*\}.
\]
\end{proposition}

\begin{proof}
The shared foreground overlap is a subspace of each individual
foreground image. Therefore its dimension cannot exceed the dimension of
either image, hence cannot exceed the minimum of the two.
\end{proof}

\begin{remark}
Proposition~\ref{prop:sharedfgdim} is only a bookkeeping result. It
does not prove that the overlap is non-trivial. Non-triviality of the
shared foreground remains a separate conditional input.
\end{remark}

\subsection{Shared stable hull}

Write
\[
\omega(\cdot)=\langle\Omega,\cdot\,\Omega\rangle
\]
on the ambient algebra \(\mathcal A_U\), and let
\(\sigma_s^{\omega,\mathrm{tot}}\) be its ambient modular flow.

\begin{definition}[Shared stable-carrier class]
\label{def:sharedcarrier}
Let \(\mathfrak F_{\mathrm{sh}}\) be a non-trivial shared foreground
overlap. Define
\[
\mathfrak L_{\mathrm{sh}}^{\mathrm{st}}
:=
\left\{
\begin{aligned}
\mathcal M\subseteq\mathcal A_U:\;&
\mathcal M \text{ is a von Neumann subalgebra},\\
&
\mathcal A_{\mathfrak F_{\mathrm{sh}}}\subseteq \mathcal M,\\
&
\sigma_s^{\omega,\mathrm{tot}}(\mathcal M)=\mathcal M
\quad \forall s\in\mathbb R,\\
&
\mathcal E_t^{(i)}(\mathcal M)\subseteq \mathcal M
\quad \forall t\ge 0,\ i=1,2
\end{aligned}
\right\}.
\]
\end{definition}

\begin{definition}[Shared stable hull]
\label{def:sharedhull}
Assume \(\mathfrak L_{\mathrm{sh}}^{\mathrm{st}}\neq \varnothing\). The
\emph{shared stable hull} of the shared foreground is
\[
\Phi_{\mathrm{sh}}^{\mathrm{st}}(\mathfrak F_{\mathrm{sh}})
:=
\bigcap_{\mathcal M\in\mathfrak L_{\mathrm{sh}}^{\mathrm{st}}}
\mathcal M.
\]
\end{definition}

\begin{proposition}[Minimal shared stable hull]
\label{prop:sharedhull}
Assume \(\mathfrak L_{\mathrm{sh}}^{\mathrm{st}}\neq\varnothing\). Then
\[
\Phi_{\mathrm{sh}}^{\mathrm{st}}(\mathfrak F_{\mathrm{sh}})
\]
is a well-defined von Neumann subalgebra of \(\mathcal A_U\) satisfying:
\begin{enumerate}
\item
  \(\mathcal A_{\mathfrak F_{\mathrm{sh}}}\subseteq
  \Phi_{\mathrm{sh}}^{\mathrm{st}}(\mathfrak F_{\mathrm{sh}})\);
\item
  \[
  \sigma_s^{\omega,\mathrm{tot}}
  \bigl(\Phi_{\mathrm{sh}}^{\mathrm{st}}(\mathfrak F_{\mathrm{sh}})\bigr)
  =
  \Phi_{\mathrm{sh}}^{\mathrm{st}}(\mathfrak F_{\mathrm{sh}})
  \qquad
  \forall s\in\mathbb R;
  \]
\item
  \[
  \mathcal E_t^{(i)}
  \bigl(\Phi_{\mathrm{sh}}^{\mathrm{st}}(\mathfrak F_{\mathrm{sh}})\bigr)
  \subseteq
  \Phi_{\mathrm{sh}}^{\mathrm{st}}(\mathfrak F_{\mathrm{sh}})
  \qquad
  \forall t\ge 0,\ i=1,2;
  \]
\item it is minimal with these properties: if
  \(\mathcal N\subseteq\mathcal A_U\) is any von Neumann subalgebra
  containing \(\mathcal A_{\mathfrak F_{\mathrm{sh}}}\) and stable
  under the ambient modular flow and both semigroups, then
  \[
  \Phi_{\mathrm{sh}}^{\mathrm{st}}(\mathfrak F_{\mathrm{sh}})
  \subseteq
  \mathcal N.
  \]
\end{enumerate}
\end{proposition}

\begin{proof}
The argument is identical to the stable-hull construction in
Foundation~I~\cite{FoundationI}, with the only change that one now
imposes stability under both source semigroups. An arbitrary
intersection of von Neumann subalgebras is again a von Neumann
subalgebra: intersections of \(\sigma\)-weakly closed
\(*\)-subalgebras remain \(\sigma\)-weakly closed \(*\)-subalgebras.
Because every member of
\(\mathfrak L_{\mathrm{sh}}^{\mathrm{st}}\) contains
\(\mathcal A_{\mathfrak F_{\mathrm{sh}}}\), so does the intersection.
Stability under the ambient modular flow and both semigroups is
preserved by intersection: if \(A\in\bigcap_{\alpha}\mathcal M_\alpha\),
then \(\mathcal E_t^{(i)}(A)\in\mathcal M_\alpha\) for every
\(\alpha\), \(t\), and \(i\), and similarly for the modular flow, hence
the image again lies in \(\bigcap_{\alpha}\mathcal M_\alpha\).
Minimality is immediate from the definition as the intersection of all
subalgebras with the required properties.
\end{proof}

\subsection{Main theorem}

\begin{theorem}[Two-source shared admissible manifest sector]
\label{thm:main}
Let \(S_{\mathrm{multi}}^{(2)}\) be a two-source datum. Assume:
\begin{enumerate}
\item \textbf{Individual source closure.} Each source separately
  satisfies the named Foundation~I hypotheses needed to produce
  \[
  (\mathcal A_c^{(1)},\omega_1),
  \qquad
  (\mathcal A_c^{(2)},\omega_2),
  \]
  together with foreground frames
  \(\mathfrak F^{(1)}\) and \(\mathfrak F^{(2)}\).
\item \textbf{Joint-stability package.} There exists a candidate shared
  sector
  \[
  \mathcal A_{\mathrm{sh}}^{\mathrm{cand}}
  \subseteq
  \mathcal A_c^{(1)}\cap\mathcal A_c^{(2)}
  \]
  satisfying JS-A, JS-B, and JS-C.
\item \textbf{Non-trivial shared foreground.} The shared foreground
  overlap \(\mathfrak F_{\mathrm{sh}}\) is non-zero and admits a
  realization by exact shared pointer-aligned observables
  \(O_1^{\mathrm{sh}},\dots,O_{k_{\mathrm{sh}}}^{\mathrm{sh}}\) lying in
  \(\mathcal A_{\mathrm{sh}}^{\mathrm{cand}}\).
\item \textbf{Admissibility input for Paper~I.} The pair
  \[
  \bigl(
  \Phi_{\mathrm{sh}}^{\mathrm{st}}(\mathfrak F_{\mathrm{sh}}),
  \omega_{\mathrm{sh}}|_{\Phi_{\mathrm{sh}}^{\mathrm{st}}(\mathfrak F_{\mathrm{sh}})}
  \bigr)
  \]
  satisfies the admissibility conditions required by
  Paper~I~\cite{PaperI}.
\end{enumerate}
Then:
\begin{enumerate}
\item the shared stable hull is well-defined and satisfies
  \[
  \Phi_{\mathrm{sh}}^{\mathrm{st}}(\mathfrak F_{\mathrm{sh}})
  \subseteq
  \mathcal A_{\mathrm{sh}}^{\mathrm{cand}}
  \subseteq
  \mathcal A_c^{(1)}\cap \mathcal A_c^{(2)};
  \]
\item the restricted state
  \[
  \omega_{\mathrm{sh}}^{\mathrm{st}}
  :=
  \omega_{\mathrm{sh}}|_{\Phi_{\mathrm{sh}}^{\mathrm{st}}(\mathfrak F_{\mathrm{sh}})}
  \]
  is faithful on
  \(\Phi_{\mathrm{sh}}^{\mathrm{st}}(\mathfrak F_{\mathrm{sh}})\);
\item the pair
  \[
  \Bigl(
  \Phi_{\mathrm{sh}}^{\mathrm{st}}(\mathfrak F_{\mathrm{sh}}),
  \omega_{\mathrm{sh}}^{\mathrm{st}}
  \Bigr)
  \]
  is a shared admissible manifest sector;
\item therefore the Paper~I inseparability theorem applies to that pair.
\end{enumerate}
\end{theorem}

\begin{proof}
By assumption~(3), the exact shared pointer-aligned observables
\(O_1^{\mathrm{sh}},\dots,O_{k_{\mathrm{sh}}}^{\mathrm{sh}}\) lie in
\(\mathcal A_{\mathrm{sh}}^{\mathrm{cand}}\). Since
\(\mathcal A_{\mathrm{sh}}^{\mathrm{cand}}\) is a von Neumann
subalgebra by Definition~2.4, it contains the generated shared
foreground algebra:
\[
\mathcal A_{\mathfrak F_{\mathrm{sh}}}
\subseteq
\mathcal A_{\mathrm{sh}}^{\mathrm{cand}}.
\]
By assumption~(2), the candidate shared sector satisfies JS-A, hence it
is invariant under the ambient modular flow and under both source
semigroups. Therefore
\[
\mathcal A_{\mathrm{sh}}^{\mathrm{cand}}
\in
\mathfrak L_{\mathrm{sh}}^{\mathrm{st}}.
\]
In particular, the class \(\mathfrak L_{\mathrm{sh}}^{\mathrm{st}}\) is
nonempty. Hence
Definition~\ref{def:sharedhull} and
Proposition~\ref{prop:sharedhull} apply, so the shared stable hull is a
well-defined von Neumann subalgebra of \(\mathcal A_U\).

Because \(\mathcal A_{\mathrm{sh}}^{\mathrm{cand}}\) is a member of the
shared stable-carrier class, minimality gives
\[
\Phi_{\mathrm{sh}}^{\mathrm{st}}(\mathfrak F_{\mathrm{sh}})
\subseteq
\mathcal A_{\mathrm{sh}}^{\mathrm{cand}}.
\]
By assumption~(2), the candidate shared sector is already contained in
\(\mathcal A_c^{(1)}\cap\mathcal A_c^{(2)}\), so
\[
\Phi_{\mathrm{sh}}^{\mathrm{st}}(\mathfrak F_{\mathrm{sh}})
\subseteq
\mathcal A_c^{(1)}\cap\mathcal A_c^{(2)}.
\]
This proves part~(1).

By JS-B, \(\omega_{\mathrm{sh}}\) is faithful on
\(\mathcal A_{\mathrm{sh}}^{\mathrm{cand}}\). The restriction of a
faithful normal state to a von Neumann subalgebra is again faithful.
Therefore
\[
\omega_{\mathrm{sh}}^{\mathrm{st}}
=
\omega_{\mathrm{sh}}|_{\Phi_{\mathrm{sh}}^{\mathrm{st}}(\mathfrak F_{\mathrm{sh}})}
\]
is faithful on the shared stable hull. This proves part~(2).

Parts~(1) and~(2), together with the stability properties from
Proposition~\ref{prop:sharedhull}, show that the pair
\[
\Bigl(
\Phi_{\mathrm{sh}}^{\mathrm{st}}(\mathfrak F_{\mathrm{sh}}),
\omega_{\mathrm{sh}}^{\mathrm{st}}
\Bigr)
\]
is a shared manifest sector in the sense intended by the present paper.
Assumption~(4) is the external admissibility input that upgrades it to a
shared \emph{admissible} manifest sector, so Paper~I applies. This
proves parts~(3) and~(4).
\end{proof}

\begin{corollary}[Shared inseparability]
\label{cor:sharedinseparability}
Under the hypotheses of Theorem~\ref{thm:main}, geometry and modular
flow on the shared sector
\[
\Bigl(
\Phi_{\mathrm{sh}}^{\mathrm{st}}(\mathfrak F_{\mathrm{sh}}),
\omega_{\mathrm{sh}}^{\mathrm{st}}
\Bigr)
\]
are not independently deformable in the sense of Paper~I. In
particular, any geometric reading
\(\mathrm{Geom}(\Phi_{\mathrm{sh}}^{\mathrm{st}}(\mathfrak F_{\mathrm{sh}}),
\omega_{\mathrm{sh}}^{\mathrm{st}})\)
defined within the Paper~I theorem domain is a \emph{shared} geometric
sector.
\end{corollary}

\begin{proof}
This is exactly the conclusion of the Paper~I theorem applied to the
pair isolated by Theorem~\ref{thm:main}.
\end{proof}

\begin{remark}[What the main theorem does \emph{not} prove]
Theorem~\ref{thm:main} does not prove a general shared-world theorem.
It proves something narrower: if a two-source datum already admits a
candidate shared sector satisfying the named compatibility and
admissibility inputs, then the overlap of the two source chains carries
an honest theorem-level manifest sector to which the Paper~I machinery
may be applied. It also does not derive the candidate shared sector from
multi-source entanglement or source dynamics alone. This is a foothold,
not the full Bridge~2 closure.
\end{remark}

\begin{remark}[Where the theorem's new content lies]
The genuine new content of Theorem~\ref{thm:main} is the conditional
inclusion chain
\[
\text{shared foreground inclusion}
\;+\;
\text{JS-A}
\Longrightarrow
\mathcal A_{\mathrm{sh}}^{\mathrm{cand}}
\in
\mathfrak L_{\mathrm{sh}}^{\mathrm{st}}
\Longrightarrow
\Phi_{\mathrm{sh}}^{\mathrm{st}}(\mathfrak F_{\mathrm{sh}})
\subseteq
\mathcal A_c^{(1)}\cap\mathcal A_c^{(2)}.
\]
The Paper~I admissibility package is not derived here. It is a named
external input specifying when the isolated shared stable hull lies in
the theorem domain of Paper~I.
\end{remark}

\section{Common/Private Decomposition}\label{sec:cp}

The main theorem isolates a shared admissible manifest sector. We now
ask how each observer's full manifest algebra relates to that shared
sector. In the Type~III\(_1\) setting, one should not treat this as a
canonical factorization problem. The correct formulation is
regime-dependent and uses a conditional expectation language.

\begin{definition}[Shared/private regime assumption]
\label{def:regime}
Let
\[
\mathcal A_{\mathrm{sh}}
\subseteq
\mathcal A_c^{(i)}
\subseteq
\mathcal A_U,
\qquad
i=1,2.
\]
We say that the pair
\((\mathcal A_c^{(i)},\mathcal A_{\mathrm{sh}})\) satisfies the
\emph{shared/private regime assumption} if for each \(i\) there exist:
\begin{enumerate}
\item a von Neumann subalgebra
  \(\mathcal A_{\mathrm{priv}}^{(i)}\subseteq \mathcal A_c^{(i)}\);
\item a faithful normal conditional expectation
  \[
  E_{\mathrm{sh}}^{(i)}:\mathcal A_c^{(i)}\to \mathcal A_{\mathrm{sh}}
  \]
  preserving the individual manifest state in the sense that
  \[
  \omega_i\circ E_{\mathrm{sh}}^{(i)}=\omega_i;
  \]
\item a generating relation
  \[
  \mathcal A_c^{(i)}
  =
  W^*\!\bigl(\mathcal A_{\mathrm{sh}}\cup
  \mathcal A_{\mathrm{priv}}^{(i)}\bigr).
  \]
\end{enumerate}
\end{definition}

\begin{remark}
Definition~\ref{def:regime} does not assert uniqueness of the private
complement, and it does not assert a tensor-product or free-product
factorization. It only says that, under a regime assumption, the full
manifest algebra can be generated from the shared sector together with a
private complement, with a faithful conditional expectation projecting
back to the shared sector.
\end{remark}

\begin{remark}[Minimal package for common/private language]
Definition~\ref{def:regime} is the minimal package needed before
common/private language can even be used honestly in the present paper.
Without a private complement, a faithful conditional expectation onto
the shared sector, and a generating relation for the full manifest
algebra, the phrase ``private part'' remains interpretive rather than
mathematical.
\end{remark}

\begin{proposition}[Shared agreement and overlap-level discrepancy]
\label{prop:cpdiscrepancy}
Assume the hypotheses of Theorem~\ref{thm:main}, and suppose the
shared/private regime assumption holds with
\[
\mathcal A_{\mathrm{sh}}
:=
\Phi_{\mathrm{sh}}^{\mathrm{st}}(\mathfrak F_{\mathrm{sh}}).
\]
Let
\[
\mathcal A_{\mathrm{ov}}
:=
\mathcal A_c^{(1)}\cap \mathcal A_c^{(2)}.
\]
Define the overlap discrepancy functional
\[
\delta_{12}
:=
\omega_1|_{\mathcal A_{\mathrm{ov}}}
-
\omega_2|_{\mathcal A_{\mathrm{ov}}}
\in
(\mathcal A_{\mathrm{ov}})_*.
\]
Then:
\begin{enumerate}
\item the two restricted manifest states agree on the shared sector:
  \[
  \omega_1|_{\mathcal A_{\mathrm{sh}}}
  =
  \omega_2|_{\mathcal A_{\mathrm{sh}}}
  =
  \omega_{\mathrm{sh}}^{\mathrm{st}}|_{\mathcal A_{\mathrm{sh}}};
  \]
\item equivalently,
  \[
  \delta_{12}|_{\mathcal A_{\mathrm{sh}}}=0;
  \]
\item if \(A\in\mathcal A_{\mathrm{ov}}\) satisfies
  \(\delta_{12}(A)\neq 0\), then \(A\notin\mathcal A_{\mathrm{sh}}\);
\end{enumerate}
\end{proposition}

\begin{proof}
For \(A\in\mathcal A_{\mathrm{sh}}\), Theorem~\ref{thm:main} gives
\[
\omega_1(A)
=
\omega_{\mathrm{sh}}^{\mathrm{st}}(A)
=
\omega_2(A).
\]
Therefore
\[
\delta_{12}(A)=0
\qquad
\forall A\in\mathcal A_{\mathrm{sh}}.
\]
This proves the first two claims. The third is the contrapositive of the
second.
\end{proof}

\begin{corollary}[Shared manifest content]
\label{cor:sharedcontent}
Under the hypotheses of Proposition~\ref{prop:cpdiscrepancy}, every
observable in the shared sector carries common manifest content.
Observer-dependent discrepancy, insofar as it is jointly accessible on
\(\mathcal A_{\mathrm{ov}}\), can only appear outside the shared
sector.
\end{corollary}

\begin{proof}
Proposition~\ref{prop:cpdiscrepancy} shows that the two restricted
manifest states agree on the shared sector and that any jointly
accessible disagreement must be detected off that sector. That is the
formal meaning of ``shared manifest content'' used in the present paper.
\end{proof}

\begin{remark}[Shared/private reading of Proposition~\ref{prop:cpdiscrepancy}]
Under the shared/private regime assumption, Proposition~\ref{prop:cpdiscrepancy}
may be read as saying that any overlap-level disagreement is detected
only outside the common manifest sector, i.e.\ on the observer-dependent
side of the decomposition language introduced in this section.
\end{remark}

\begin{remark}[What Section~\ref{sec:cp} does \emph{not} prove]
Section~\ref{sec:cp} does not prove existence of a canonical private
factor, nor does it derive the shared/private regime assumption from the
bare two-source datum. Its role is narrower: once a shared manifest
sector has been isolated and a regime assumption is granted, it makes
precise where observer agreement must hold and where observer-dependent
discrepancy may still appear.
\end{remark}

\begin{remark}[No canonical factorization claim]
Corollary~\ref{cor:sharedcontent} should not be read as a theorem that
\[
\mathcal A_c^{(i)}
\cong
\mathcal A_{\mathrm{sh}}\bar\otimes
\mathcal A_{\mathrm{priv}}^{(i)}
\]
or any analogous canonical factorization holds. The paper makes no such
claim. It uses only the weaker and safer statement that, under a regime
assumption, a shared sector and a private complement jointly generate
the full manifest algebra.
\end{remark}

\section{Conditional \texorpdfstring{$N$}{N}-Source Extension}\label{sec:nsource}

The two-source theorem is the paper's main result. The natural next step
is a conditional \(N\)-source extension. This section records the
generalization in the same honest style as Foundation~I: whatever can be
internalized immediately is internalized, and the genuinely new input is
kept explicit.

\begin{definition}[Auxiliary \(N\)-source conditions]
\label{def:Nconds}
Let
\[
S_{\mathrm{multi}}
=
(\mathcal H_U,\mathcal A_U,\Omega,\{\mathcal E_t^{(i)}\}_{i=1}^N)
\]
be a multi-source datum with individual manifest outputs
\((\mathcal A_c^{(i)},\omega_i)_{i=1}^N\). We use the following
auxiliary conditions.
\begin{enumerate}
\item \textbf{Exchangeability:} the source family is invariant, at the
  level relevant to the theorem, under finite permutations of the
  labels.
\item \textbf{Pairwise threshold condition:} for every pair
  \(i\neq j\), the two-source hypotheses of
  Theorem~\ref{thm:main} hold relative to the same candidate shared
  sector \(\mathcal A_{\mathrm{sh}}^{\mathrm{cand}}\).
\item \textbf{Graph-compatibility condition:} the source labels are
  vertices of a connected graph \(G\), and every edge
  \((i,j)\in E(G)\) satisfies the two-source hypotheses relative to the
  same candidate shared sector.
\item \textbf{Global consensus functional:} there exists a functional
  \(C_N\) whose only role in the present paper is to define the domain
  on which the previous conditions are asserted. No quantitative bound
  involving \(C_N\) is derived here.
\end{enumerate}
\end{definition}

\begin{theorem}[Conditional \(N\)-source shared manifest sector]
\label{thm:N}
Let \(S_{\mathrm{multi}}\) be an \(N\)-source datum with individual
manifest outputs \((\mathcal A_c^{(i)},\omega_i)_{i=1}^N\). Assume:
\begin{enumerate}
\item there exists a candidate shared sector
  \[
  \mathcal A_{\mathrm{sh}}^{\mathrm{cand}}
  \subseteq
  \bigcap_{i=1}^N \mathcal A_c^{(i)}
  \]
  satisfying JS-A and JS-B for all sources;
\item the pairwise-threshold condition of
  Definition~\ref{def:Nconds} holds, with a compatible family of shared
  foreground overlaps generating a common foreground algebra
  \(\mathcal A_{\mathfrak F_{\mathrm{sh}}}\);
\item the \(N\)-source stable-carrier class
  \[
  \mathfrak L_{\mathrm{sh},N}^{\mathrm{st}}
  :=
  \left\{
  \mathcal M\subseteq\mathcal A_U:
  \begin{array}{l}
  \mathcal M \text{ is a von Neumann subalgebra},\\
  \mathcal A_{\mathfrak F_{\mathrm{sh}}}\subseteq \mathcal M,\\
  \sigma_s^{\omega,\mathrm{tot}}(\mathcal M)=\mathcal M,\ \forall s,\\
  \mathcal E_t^{(i)}(\mathcal M)\subseteq\mathcal M,\ \forall t,\forall i
  \end{array}
  \right\}
  \]
  is nonempty and contains
  \(\mathcal A_{\mathrm{sh}}^{\mathrm{cand}}\);
\item the resulting \(N\)-source shared stable hull satisfies the
  admissibility hypotheses needed by Paper~I.
\end{enumerate}
Then the \(N\)-source shared stable hull is a shared admissible manifest
sector contained in
\(\bigcap_{i=1}^N\mathcal A_c^{(i)}\).
\end{theorem}

\begin{proof}
The proof is the same minimality argument as in the two-source case,
with the only change that one now intersects over all source labels
rather than over two. The candidate shared sector is, by assumption, a
member of the \(N\)-source stable-carrier class. Therefore the
intersection defining the \(N\)-source shared stable hull is
well-defined and lies inside the candidate sector. The remaining steps
repeat the proof of Theorem~\ref{thm:main}.
\end{proof}

\begin{remark}[What Theorem~\ref{thm:N} does not contain]
Theorem~\ref{thm:N} does not provide:
\begin{enumerate}
\item a uniqueness or maximality theorem for the \(N\)-source shared
  sector;
\item a quantitative relationship between the number/strength of sources
  and geometric stability;
\item a proof that the auxiliary conditions of
  Definition~\ref{def:Nconds} follow from the bare multi-source datum.
\end{enumerate}
All three belong to later work.
\end{remark}

\section{Open Problems}\label{sec:open}

The conditional inputs of Foundation~II are not cosmetic. They isolate
the genuinely unfinished parts of the programme.

Heuristically, the present two-source theorem is closer to a
mean-field foothold than to a full correlated treatment: it isolates a
shared effective manifest sector in a restricted regime, while genuinely
cross-source structure remains the central open problem for any future
multi-source closure.

\begin{enumerate}
\item \textbf{Model-independent joint stability.} The present paper
  names JS-A, JS-B, and JS-C. It does not derive them model-independently
  from a bare multi-source datum. The analogue of Foundation~I's
  observer-emergence bottleneck therefore persists in multi-source form.

\item \textbf{Compatible Nakajima--Zwanzig projections.} The present
  notion of compatible persistent projections is a new definition, not a
  standard theorem. A general existence theory for such compatible
  projections remains open.

\item \textbf{Shared-foreground non-triviality.} The paper proves that
  if a non-trivial shared foreground overlap exists, then its dimension
  is bounded by the individual budget ceilings. It does not prove that
  the overlap must be non-zero.

\item \textbf{Quantitative consensus versus stability.} The paper
  deliberately avoids a bound of the form
  \[
  \text{consensus strength}
  \Longrightarrow
  \text{geometric fluctuation bound}.
  \]
  That is the separate Bridge~3 problem.

\item \textbf{\(N\)-source non-conditional closure.} The present
  \(N\)-source theorem is conditional. A truly non-conditional closure
  would require deriving the pairwise-threshold, graph, or
  exchangeability inputs from source-side dynamics.

\item \textbf{Uniqueness and maximality of the shared sector.} Even in
  the two-source case, the present theorem proves only existence of a
  minimal shared stable hull inside a candidate shared carrier. It does
  not prove that this sector is unique among all plausible shared
  manifest sectors, or maximal in any global sense.

\item \textbf{Relation to Paper~II.} Paper~II constrains higher-order
  geometry by multi-directional consistency~\cite{PaperII}. The present
  paper does not yet prove how the shared manifest sector of
  Foundation~II interacts with the Paper~II small-ball remainder
  structure. That interface remains open.

\item \textbf{Type~III and infinite-dimensional obstacles.} The
  common/private regime assumption is not automatic in the Type~III
  setting. Conditional expectations, complements, and generating
  decompositions all become more delicate in infinite dimensions. In
  particular, existence of a faithful normal conditional expectation
  onto a candidate shared sector typically interacts with modular
  invariance in the sense of Takesaki and may fail outside suitable
  invariant regimes. Even when such expectations exist, complements and
  generating relations need not be unique, and non-factorial or
  center-bearing sectors may force a central decomposition before any
  shared/private language becomes well behaved.
\end{enumerate}

\section{Conditional Input Ledger}\label{sec:ledger}

This section records the named inputs that remain conditional in the
present paper.

\begin{enumerate}
\item \textbf{I1: Individual source closure.} Each source separately
  satisfies the named Foundation~I hypotheses needed to produce an
  individual manifest pair \((\mathcal A_c^{(i)},\omega_i)\).

\item \textbf{JS-A: Shared-sector invariance and semigroup
  compatibility.} The candidate shared sector is stable under all
  relevant semigroups and under the ambient modular flow, and the
  restricted semigroups commute there.

\item \textbf{JS-B: Shared-state existence and faithfulness.} A common
  faithful normal state exists on the candidate shared sector and is
  compatible with the source-specific manifest states.

\item \textbf{JS-C: Compatible persistent projections.} The pair of
  source-side Nakajima--Zwanzig projections commute and admit a common
  projected kernel on the shared sector.

\item \textbf{F1: Non-trivial shared foreground with exact pointer
  realization.} The overlap of the two source-side foreground
  realizations is non-zero, finite dimensional, and admits an exact
  pointer-aligned realization inside the candidate shared sector.

\item \textbf{A1: Admissibility input for Paper~I.} The shared stable
  hull satisfies the admissibility package needed for the
  geometry--modular-flow inseparability theorem.

\item \textbf{SP1: Shared/private regime assumption.} Each individual
  manifest algebra admits a faithful conditional expectation onto the
  shared sector together with a private complement that generates the
  full algebra jointly with the shared sector.

\item \textbf{N1: \(N\)-source auxiliary conditions.} Exchangeability,
  pairwise threshold, graph compatibility, or a global consensus
  functional are available when one moves beyond the two-source case.
\end{enumerate}

\appendix

\section{Imported Interfaces}
\label{app:premises}

This appendix records, in theorem-ready form, the external interfaces
imported from Genesis, Foundation~I, and the Consciousness Trilogy. They
are not reproved here, but they are restated explicitly so that the
present paper can be read without reconstructing its dependencies from
the earlier series papers. What is \emph{not} imported should also be
kept clear: the multi-source notions of candidate shared sector, JS-A,
JS-B, JS-C, shared stable hull, and shared/private regime are defined
internally in Foundation~II itself.

\begin{proposition}[Foundation~I single-source closure package]
\label{prop:appF1}
Foundation~I~\cite{FoundationI} isolates a conditional chain
\[
S_{\mathrm{source}}
\to
\mathcal O_A
\to
\mathcal O_B
\to
(\mathcal A_c,\omega),
\]
where:
\begin{enumerate}
\item Stage~A supplies a projected record object
  \[
  \mathcal O_A=(\mathcal P,\omega_t,\mathcal K_*(t,s)),
  \]
  where \(\mathcal P:\mathcal A_U\to\mathcal A_U\) is a normal
  idempotent projection and \(\mathcal K_*(t,s)\) is a dual
  non-Markovian memory kernel;
\item Stage~B supplies a foreground frame \(\mathfrak F\) with a
  finite-dimensional foreground space
  \(V_{\mathrm{fg}}(\mathfrak F)\) and an operator realization
  \[
  \iota:V_{\mathrm{fg}}(\mathfrak F)\hookrightarrow \mathcal A_U^{\mathrm{sa}};
  \]
\item Stage~C supplies a manifest pair
  \[
  (\mathcal A_c,\omega)
  =
  \bigl(\Phi_{\mathfrak F}^{\mathrm{st}}(\mathfrak F),\omega\bigr),
  \]
  where \(\Phi_{\mathfrak F}^{\mathrm{st}}(\mathfrak F)\subseteq
  \mathcal A_U\) is a minimal modular-decoherence stable hull, hence a
  von Neumann subalgebra carrying a faithful normal restricted state;
\item in the exact fixed-point comparison regime imported from
  Foundation~I, one has the inclusion
  \[
  \Phi_{\mathfrak F}^{\mathrm{st}}(\mathfrak F)
  \subseteq
  \mathcal A_c^{\mathrm{HAFF}}.
  \]
  (not used in the present paper; recorded for later reference).
\end{enumerate}
\end{proposition}

\begin{proposition}[Imported single-source interfaces used in Foundation~II]
\label{prop:appinterfaces}
Whenever Foundation~II invokes the \(i\)-th source chain, the only
external single-source objects it uses are:
\begin{enumerate}
\item a Stage~A projection
  \(\mathcal P^{(i)}\), projected normal states
  \(\omega_t^{(i)}\), and dual memory kernel
  \(\mathcal K_{*,i}(t,s)\);
\item a Stage~B foreground frame \(\mathfrak F^{(i)}\) with realized
  finite-dimensional foreground image
  \[
  \iota_i\bigl(V_{\mathrm{fg}}(\mathfrak F^{(i)})\bigr)
  \subseteq
  \mathcal A_U^{\mathrm{sa}};
  \]
\item a Stage~C manifest pair
  \[
  (\mathcal A_c^{(i)},\omega_i)
  =
  \bigl(\Phi_{\mathfrak F^{(i)}}^{\mathrm{st}}(\mathfrak F^{(i)}),
  \omega_i\bigr).
  \]
\end{enumerate}
No multi-source compatibility theorem is imported at this stage.
Foundation~II introduces the multi-source compatibility layer itself.
\end{proposition}

\begin{proposition}[Paper~I manifest inseparability theorem]
\label{prop:appP1}
Paper~I~\cite{PaperI} proves that for an admissible manifest pair
\((\mathcal A_c,\omega)\), geometry and modular flow are not
independently deformable. In particular, once a pair
\((\mathcal A_c,\omega)\) in the theorem domain is given, the geometry
and modular sectors should be treated as two constrained readings of one
underlying algebraic datum. Foundation~II uses this result only through
its admissibility interface: whenever a shared stable hull is shown, by
named external input, to belong to the Paper~I theorem domain, the
inseparability conclusion becomes available.
\end{proposition}

\begin{proposition}[Paper~II as a consistency constraint]
\label{prop:appP2}
Paper~II~\cite{PaperII} shows that in the CHM small-ball setting,
multi-directional consistency of geometry and modular data forces the
linearized Einstein equation. The present paper uses this only as a
boundary calibration: it does not import the Paper~II theorem as a
formal input to the shared-sector existence proof.
\end{proposition}

\begin{proposition}[T-DOME source-side ingredients]
\label{prop:appTDOME}
The T-DOME papers provide the three single-source ingredients used
upstream of the present paper:
\begin{enumerate}
\item T-DOME~I: persistent far-from-equilibrium subsystems require
  non-Markovian memory and thereby motivate the Stage~A projected record
  process with kernel \(\mathcal K_*(t,s)\)~\cite{TDOMEI,BreuerPetruccione2002};
\item T-DOME~II: finite processing budgets force foreground/background
  partition, symmetry-broken frame selection, and a finite effective
  foreground dimension \(k^*\)~\cite{TDOMEII};
\item T-DOME~III: self-calibration enters through a Fisher-information
  prior carried by the foreground structure, but no direct theorem from
  T-DOME~III is used in the operator-algebraic proofs below beyond this
  source-side calibration role~\cite{TDOMEIII}.
\end{enumerate}
\end{proposition}

\begin{proposition}[HAFF ambient setting]
\label{prop:appHAFF}
The HAFF programme supplies the ambient structural setting used here:
\begin{enumerate}
\item geometry is read from accessible algebraic structure rather than
  supplied as a background container~\cite{HAFFA};
\item accessible algebras are selected by stability under the relevant
  dynamics rather than by subjective choice~\cite{HAFFD};
\item in the exact fixed-point regime used for Foundation~I and imported
  here, the HAFF accessible algebra
  \[
  \mathcal A_c^{\mathrm{HAFF}}\subseteq \mathcal A_U
  \]
  is a von Neumann subalgebra stable under the relevant semigroup and
  the ambient modular flow;
\item fixed-point and stability ideas for completely positive semigroups
  motivate the use of stable carrier classes and exact pointer sectors
  in restricted regimes~\cite{FrigerioVerri1982,Zurek2003}.
\end{enumerate}
\end{proposition}

\begin{remark}[What earlier papers do and do not supply]
The earlier Genesis and Consciousness papers supply only single-source
objects and single-source comparison statements of the kind listed above.
They do \emph{not} supply a multi-source theorem asserting that two
projections commute, that two foreground frames overlap nontrivially, or
that two source chains admit a common stationary state. Those are the
precise new conditional interfaces named JS-A, JS-B, and JS-C in the
main text.
\end{remark}

\begin{remark}[Conditional expectations]
The common/private regime assumption uses faithful normal conditional
expectations in the sense of Takesaki~\cite{Takesaki1972}. No claim is
made here that such expectations always exist onto the shared sector in
the Type~III\(_1\) setting; their existence is part of the regime
assumption itself.
\end{remark}
\newpage
\addcontentsline{toc}{section}{References}

\begin{thebibliography}{99}
\sloppy

\bibitem{FoundationI}
S.~Liu,
\emph{Foundation I: From Source to Manifestation: A Conditional Same-Root Program for Geometry and Modular Flow},
Consciousness Series (2026),
Zenodo, DOI:~10.5281/zenodo.19627786.

\bibitem{PaperI}
S.~Liu,
\emph{The Inseparability of Geometry and Modular Flow: A Structural Inseparability Theorem and Its Implications for the Hard Problem},
Consciousness Trilogy, Paper~I (2026),
Zenodo, DOI:~10.5281/zenodo.19607387.

\bibitem{PaperII}
S.~Liu,
\emph{Higher-Dimensional Consistency of Geometry and Modular Flow: Small-Ball Modular Data and Linearized Gravity},
Consciousness Trilogy, Paper~II (2026),
Zenodo, DOI:~10.5281/zenodo.19607543.

\bibitem{PaperIII}
S.~Liu,
\emph{The Structural Horizon of Self-Description: Fisher Bounds, Small-Ball Remainders, and the Quantitative Explanatory Gap},
Consciousness Trilogy, Paper~III (2026),
Zenodo, DOI:~10.5281/zenodo.19607724.

\bibitem{HAFFA}
S.~Liu,
\emph{Emergent Geometry from Coarse-Grained Observable Algebras: The Holographic Alaya-Field Framework},
Zenodo (2026), DOI:~10.5281/zenodo.18361706.

\bibitem{HAFFD}
S.~Liu,
\emph{Gravitational Phenomena as Emergent Properties of Observable Algebra Selection: A Structural Analysis},
Zenodo (2026), DOI:~10.5281/zenodo.18388881.

\bibitem{TDOMEI}
S.~Liu,
\emph{Non-Markovian Memory and the Thermodynamic Necessity of Temporal Accumulation},
Zenodo (2026), DOI:~10.5281/zenodo.18574342.

\bibitem{TDOMEII}
S.~Liu,
\emph{Spontaneous Symmetry Breaking of Reference Frames as a Computational Cost Minimization Strategy},
Zenodo (2026), DOI:~10.5281/zenodo.18579703.

\bibitem{TDOMEIII}
S.~Liu,
\emph{Fisher Information Geometry and the Thermodynamic Cost of Self-Referential Calibration},
Zenodo (2026), DOI:~10.5281/zenodo.18591771.

\bibitem{BreuerPetruccione2002}
H.-P.~Breuer and F.~Petruccione,
\emph{The Theory of Open Quantum Systems},
Oxford University Press (2002).

\bibitem{Nakajima1958}
S.~Nakajima,
\emph{On quantum theory of transport phenomena},
Prog.\ Theor.\ Phys.\ \textbf{20}, 948--959 (1958), DOI:~10.1143/PTP.20.948.

\bibitem{Zwanzig1960}
R.~Zwanzig,
\emph{Ensemble method in the theory of irreversibility},
J.\ Chem.\ Phys.\ \textbf{33}, 1338--1341 (1960), DOI:~10.1063/1.1731409.

\bibitem{Takesaki1972}
M.~Takesaki,
\emph{Conditional expectations in von Neumann algebras},
J.\ Funct.\ Anal.\ \textbf{9}, 306--321 (1972).

\bibitem{FrigerioVerri1982}
A.~Frigerio and M.~Verri,
\emph{Long-time asymptotic properties of dynamical semigroups on \(W^*\)-algebras},
Math.\ Z.\ \textbf{180}, 275--286 (1982), DOI:~10.1007/BF01318911.

\bibitem{Zurek2003}
W.~H.~Zurek,
\emph{Decoherence, einselection, and the quantum origins of the classical},
Rev.\ Mod.\ Phys.\ \textbf{75}, 715--775 (2003), DOI:~10.1103/RevModPhys.75.715.

\bibitem{Zurek2009}
W.~H.~Zurek,
\emph{Quantum Darwinism},
Nature Phys.\ \textbf{5}, 181--188 (2009). DOI:~10.1038/nphys1202.

\bibitem{BrandaoPianiHorodecki2015}
F.~G.~S.~L.~Brand\~ao, M.~Piani, and P.~Horodecki,
\emph{Generic emergence of classical features in quantum Darwinism},
Nat.\ Commun.\ \textbf{6}, 7908 (2015), DOI:~10.1038/ncomms8908.

\bibitem{KorbiczHorodeckiHorodecki2014}
J.~K.~Korbicz, P.~Horodecki, and R.~Horodecki,
\emph{Objectivity in a noisy photonic environment through quantum state information broadcasting},
Phys.\ Rev.\ Lett.\ \textbf{112}, 120402 (2014), DOI:~10.1103/PhysRevLett.112.120402.

\end{thebibliography}
\end{document}
